% -------------------------------------------------------------------------
% WBS ID: RES-DAT-CAS
% Deliverable: Integrated Tsunami Case-Study Specifications
% Status: In progress
% Evidence status: Integrated Tohoku data-source and corridor specification established
% Review status: Not reviewed
% Last updated: 2026-07-19
% Dependencies: RES-DAT-GEO, RES-DAT-SRC, RES-DAT-OBS, RES-DAT-SCN, RES-DAT-UNC, RES-GEO-SPEC
% Downstream interfaces: RES-VER-SPEC, Software data handling, local 3D replay boundary
% -------------------------------------------------------------------------

\section{RES-DAT-CAS --- Integrated Tsunami Case-Study Specifications}
\label{sec:res-dat-cas}

\subsection{Purpose and Scope}
\label{sec:res-dat-cas-purpose}

This Deliverable replaces the previous mandatory primary-plus-secondary-event
specification with one authoritative event-reconstruction case, one validated
event baseline, bounded impact-testing variants, and documented
defence-specific local cases.

The active integrated case is the 2011 Tōhoku event with two
regional-to-local corridors: Kamaishi and Sendai coastal plain. The
case supports impact modelling, local comparison and barrier-class
comparison. ML, surrogate optimisation and generative geometry are
deferred until the hydrodynamic data pipeline and validation hierarchy
are accepted.

\subsection{Research Questions}
\label{sec:res-dat-cas-research-questions}

The integrated case-study decision asks:

\begin{enumerate}
  \item Which data sources define the event, terrain, observations,
    corridors and optional sensitivities?
  \item Which source is selected, fallback, comparison-only,
    optional or missing?
  \item Which downstream deliverable and software asset consumes each
    source?
  \item Which limitations prevent a source from being used as
    validation or solver input?
\end{enumerate}

\subsection{Terminology and Definitions}
\label{sec:res-dat-cas-definitions}

% TODO: Define the technical terminology, symbols, quantities and
% classifications used in this Deliverable.

\subsection{Literature Evidence}
\label{sec:res-dat-cas-literature-evidence}

% TODO: Record evidence from peer-reviewed papers, authoritative datasets,
% standards, technical reports and primary documentation.
%
% Do not create a detached literature-summary list. Explain how each source
% supports, contradicts or limits a project decision.

The integrated case combines official data products and peer-reviewed
Tohoku evidence. Event identity is anchored to the USGS event record
\parencite{usgs2011tohokuEvent}. Candidate finite-fault and tsunami
source products are retained from USGS finite-fault documentation,
Hayes rapid source characterisation, Fujii--Satake and DART-constrained
hindcast evidence
\parencite{usgsFiniteFaults,hayes2011rapid,fujii2011tsunamisource,wei2013japan}.
Terrain sources are drawn from J-EGG500, GSI DEM, GEBCO\_2026 and ETOPO
2022 \parencite{jodcJegg500,gsiDemDownload,gebco2026grid,noaa2022etopo}.
Observation sources include NOAA/NCEI DART, the NCEI/WDS historical
tsunami database, Mori/JTS and published observation syntheses
\parencite{noaa2005dart,noaa2011japanDart,noaaNceiHistoricalTsunami,mori2012nationwide,fine2013japan}.
ERA5 and JRA-55 are retained only as optional wind/reanalysis
sensitivity sources, not baseline hydrodynamic forcing
\parencite{copernicusEra5,jmaJra55}.

\subsection{Governing Relations and Quantitative Definitions}
\label{sec:res-dat-cas-governing-relations}

% TODO: Record relevant equations, dimensionless groups, empirical models,
% extraction rules or quantitative definitions.
%
% Use align environments for displayed equations.
% Every displayed equation must have a unique \label{eq:...}.
% Do not place punctuation after displayed equations.

\subsection{Required Data Variables and Units}
\label{sec:res-dat-cas-required-data-variables-units}

The integrated case-study manifest shall include:

\begin{itemize}
  \item event metadata, finite-fault source version and derived bed
    increments;
  \item terrain source tiles and derived mesh-projected bed fields;
  \item observation bundles classified as calibration, validation or
    diagnostic;
  \item corridor polygons for $T_{\mathrm{K}}$ and $T_{\mathrm{S}}$;
  \item local barrier-class definitions and terrain/geometry patches;
  \item replay-boundary histories for local 3D inlet forcing;
  \item provenance records \(\mathcal{P}_d\) for every source and
    derived product.
\end{itemize}

\subsection{Candidate Data Sources}
\label{sec:res-dat-cas-candidate-data-sources}

\begin{landscape}
\scriptsize
\setlength{\tabcolsep}{3pt}
\begin{longtable}{
    p{0.09\linewidth}
    p{0.08\linewidth}
    p{0.07\linewidth}
    p{0.10\linewidth}
    p{0.17\linewidth}
    p{0.10\linewidth}
    p{0.14\linewidth}
    p{0.10\linewidth}
  }
  \caption{RES-DAT integrated data-source decision table for the Tōhoku corridor case}
  \label{tab:res-dat-cas-data-source-decision-table}
  \\
  \hline
  Dataset/source name
  &
  Provider
  &
  Citation key
  &
  URL
  &
  Coverage, resolution, CRS and vertical datum
  &
  Licence/usage note
  &
  Project role and downstream deliverable
  &
  Preprocessing required and limitation
  \\
  \hline
  \endfirsthead
  \hline
  Dataset/source name
  &
  Provider
  &
  Citation key
  &
  URL
  &
  Coverage, resolution, CRS and vertical datum
  &
  Licence/usage note
  &
  Project role and downstream deliverable
  &
  Preprocessing required and limitation
  \\
  \hline
  \endhead
  J-EGG500
  &
  Japan Oceanographic Data Center
  &
  \texttt{jodcJegg500}
  &
  \url{https://www.jodc.go.jp/jodcweb/JDOSS/infoJEGG.html}
  &
  Spatial: Japan surrounding seas. Temporal: static compiled bathymetry with heterogeneous survey dates. Resolution: 500 m grid. CRS/datum: confirm from delivered metadata before ingest; bathymetric sign convention requires conversion to positive-upward $b$.
  &
  Usage/licence to be confirmed from JODC terms before redistribution.
  &
  Preferred Japan offshore bathymetry candidate for `RES-DAT-GEO`, `RES-DAT-CAS`, regional mesh and corridor terrain.
  &
  Reproject, convert depth sign/datum, clip corridors, interpolate to meshes. Limitation: smoothing and mixed data quality may suppress harbour-scale relief.
  \\
  GEBCO\_2026 Grid
  &
  GEBCO Bathymetric Compilation Group
  &
  \texttt{gebco2026grid}
  &
  \url{https://www.gebco.net/data-products/gridded-bathymetry-data}
  &
  Spatial: global ocean/land terrain. Temporal: 2026 release. Resolution: 15 arc-second grid. CRS: geographic grid. Vertical datum: treated as mean sea level, with shallow-water datum caveat.
  &
  Public-domain use with required attribution/disclaimer.
  &
  Selected outer-domain and first fallback topobathymetry for `RES-DAT-GEO`, regional domain and source-to-coast propagation.
  &
  Clip outer domain, compare with J-EGG500/GSI boundaries, resample to mesh. Limitation: not sufficient for local defence geometry or high-resolution inundation features.
  \\
  NOAA ETOPO 2022
  &
  NOAA/NCEI
  &
  \texttt{noaa2022etopo}
  &
  \url{https://www.ncei.noaa.gov/products/etopo-global-relief-model}
  &
  Spatial: global relief. Temporal: 2022 release. Resolution: 15, 30 and 60 arc-second products. CRS/datum: confirm product metadata; elevation/bathymetry grid for global relief.
  &
  NOAA citation and DOI required; redistribution terms to be checked with source package.
  &
  Secondary global fallback and audit source for `RES-DAT-GEO` and outer-domain sensitivity.
  &
  Clip/resample and compare against GEBCO. Limitation: fallback only for local corridors and not baseline local terrain.
  \\
  GSI DEM
  &
  Geospatial Information Authority of Japan
  &
  \texttt{gsiDemDownload}
  &
  \url{https://service.gsi.go.jp/kiban/app/help/}
  &
  Spatial: Japan onshore. Temporal: source-age varies by DEM class. Resolution: 1 m, 5 m and 10 m mesh products. CRS/datum: confirm in JPGIS/GML metadata before ingest.
  &
  GSI usage terms and attribution required; redistribution status to be confirmed.
  &
  Preferred onshore topography for `RES-DAT-GEO`, run-up/inundation comparison and local terrain patches.
  &
  Convert JPGIS/GML, reproject, vertical-datum reconcile, clip corridors. Limitation: mixed DEM classes and coverage gaps require provenance flags.
  \\
  USGS Tōhoku event record
  &
  U.S. Geological Survey
  &
  \texttt{usgs2011tohokuEvent}
  &
  \url{https://earthquake.usgs.gov/earthquakes/eventpage/official20110311054624120_30/executive}
  &
  Spatial: epicentre near $38.297^{\circ}\mathrm{N}$, $142.373^{\circ}\mathrm{E}$. Temporal: 2011-03-11 05:46:24 UTC. Resolution: event metadata. CRS/datum: geographic coordinates, depth metadata.
  &
  USGS public information; cite source record.
  &
  Event identity and corridor origin for `RES-DAT-SRC` and `RES-DAT-SCN`.
  &
  Store immutable event metadata and convert to projected corridor origin. Limitation: does not define seabed displacement.
  \\
  USGS finite-fault family and Hayes rapid source
  &
  U.S. Geological Survey
  &
  \texttt{usgsFiniteFaults}; \texttt{hayes2011rapid}
  &
  \url{https://earthquake.usgs.gov/data/finitefault/}
  &
  Spatial: rupture/subfault footprint. Temporal: event rupture history. Resolution: subfault dependent. CRS/datum: source-model metadata required.
  &
  Cite USGS record and article; data-package usage terms to be checked.
  &
  Candidate source baseline for `RES-DAT-SRC` and moving-bed forcing.
  &
  Parse subfaults, project geometry, compute $\Delta b_m$, rasterise to regional mesh. Limitation: calibration observations must not be reused as independent validation.
  \\
  Fujii--Satake Tōhoku source
  &
  Fujii et al.
  &
  \texttt{fujii2011tsunamisource}
  &
  DOI in bibliography
  &
  Spatial: 2011 Tōhoku source region. Temporal: event source reconstruction. Resolution: source-model dependent. CRS/datum: paper/model metadata required.
  &
  Journal citation required.
  &
  Comparison source and provenance anchor for `RES-DAT-SRC` and observation de-tiding dependence.
  &
  Reconstruct source-to-bed projection if selected. Limitation: survey corrections in Mori et al. partly depend on FS40v46-style modelling.
  \\
  NOAA/NCEI DART ocean bottom pressure
  &
  NOAA/NCEI
  &
  \texttt{noaa2005dart}; \texttt{noaa2011japanDart}
  &
  \url{https://www.ncei.noaa.gov/products/natural-hazards/tsunamis-earthquakes-volcanoes/tsunamis/dart-ocean-bottom-pressure}
  &
  Spatial: DART station points, Pacific/global network. Temporal: event time series and archive records. Resolution: includes 15-second bottom-pressure records where available. CRS/datum: station metadata and pressure-to-water-level conversion required.
  &
  Cite NOAA/NCEI; data access unrestricted through archive tools unless station-specific terms state otherwise.
  &
  Offshore arrival, phase, amplitude and source-calibration audit for `RES-DAT-OBS` and `RES-VER-SPEC`.
  &
  Download station records, de-tide/filter, classify calibration/validation. Limitation: DART records used for source inference are calibration, not independent validation.
  \\
  NCEI/WDS Global Historical Tsunami Database
  &
  NOAA/NCEI/WDS
  &
  \texttt{noaaNceiHistoricalTsunami}
  &
  \url{https://www.ncei.noaa.gov/products/natural-hazards/tsunamis-earthquakes-volcanoes/tsunamis/global-historical-data}
  &
  Spatial: global tsunami source and run-up records. Temporal: historical-to-present database. Resolution: event/run-up record level. CRS/datum: per-record metadata; datum may be ambiguous.
  &
  DOI citation required; access unrestricted according to metadata.
  &
  Observation discovery, arrival-time and max-water-height screening for `RES-DAT-OBS`.
  &
  Query Tōhoku records, cross-check against field survey. Limitation: heterogeneous historical provenance; not sole validation source.
  \\
  Mori/JTS field survey
  &
  Mori, Takahashi and 2011 Tohoku Earthquake Tsunami Joint Survey Group
  &
  \texttt{mori2012nationwide}
  &
  DOI in bibliography
  &
  Spatial: approximately 2,000 km of Japanese coastline. Temporal: post-event 2011 survey. Resolution: surveyed point/extent records. CRS/datum: survey metadata and tide corrections required.
  &
  Journal citation required; raw-data access/licence to be confirmed.
  &
  Principal run-up, inundation-height and inundation-distance evidence for `RES-DAT-OBS`, Sendai and Sanriku/Kamaishi corridor validation.
  &
  Filter by corridor, datum-normalise, classify quality. Limitation: primarily maximum water-level and extent data, not velocity or pressure histories.
  \\
  Fine et al. observation/modelling synthesis
  &
  Fine et al.
  &
  \texttt{fine2013japan}
  &
  DOI in bibliography
  &
  Spatial: Japan 2011 tsunami observation/model context. Temporal: 2011 event. Resolution: paper-level synthesis. CRS/datum: underlying records must be obtained separately.
  &
  Journal citation required.
  &
  Supports observation hierarchy and regional propagation interpretation for `RES-DAT-OBS`.
  &
  Use as evidence, not raw data. Limitation: cannot replace station/survey archives.
  \\
  ERA5 and JRA-55
  &
  Copernicus Climate Change Service; Japan Meteorological Agency
  &
  \texttt{copernicusEra5}; \texttt{jmaJra55}
  &
  \url{https://climate.copernicus.eu/climate-reanalysis}; \url{https://www.data.jma.go.jp/jra/html/JRA-55/index_en.html}
  &
  Spatial: global atmospheric reanalysis. Temporal: multi-decadal products. Resolution: product dependent. CRS/datum: atmospheric grids.
  &
  Product-specific acknowledgement and terms required.
  &
  Optional wind/reanalysis sensitivity only; not baseline hydrodynamic forcing.
  &
  No baseline preprocessing. Limitation: not used unless a sensitivity question is opened and justified.
  \\
  Kamaishi bay-mouth breakwater defence records
  &
  TODO--provider
  &
  TODO--missing-source
  &
  TODO
  &
  Spatial: Kamaishi Bay and defence footprint. Temporal: pre-2011, 2011 damage/response and reconstruction states. Resolution/CRS/datum: unresolved.
  &
  Licence unresolved.
  &
  Required for real defence recreation and wall-type reference geometry in `RES-GEO` and local `RES-MOD`/`RES-VER`.
  &
  Locate authoritative drawings, surveys, damage reports or peer-reviewed defence papers. Limitation: no citation is created until source metadata are confirmed.
  \\
  \hline
\end{longtable}
\end{landscape}

\subsection{Spatial and Temporal Coverage}
\label{sec:res-dat-cas-spatial-temporal-coverage}

The integrated case covers the 2011 Tōhoku source-to-coast event and
two active regional-to-local corridor bundles:

\begin{enumerate}
  \item $T_{\mathrm{K}}$: epicentre to Kamaishi region, with
    Sanriku/ria, bay and defence-interaction emphasis;
  \item $T_{\mathrm{S}}$: epicentre to Sendai coastal plain, with
    coastal-plain inundation emphasis.
\end{enumerate}

The local 3D replay interval shall include all waves that materially
affect run-up, overtopping, transmitted/reflected energy and load
metrics, not only the first crest.

\subsection{Coordinate and Datum Requirements}
\label{sec:res-dat-cas-coordinate-datum-requirements}

All source coordinates are acquired as longitude/latitude and then
transformed into the selected projected metric CRS for regional and
local modelling. Vertical data shall be reconciled into the
positive-upward bed convention used by `RES-MOD` and `RES-NUM`.
Observation heights shall preserve their source datum and tide
correction before any comparison to simulated $\eta$, run-up or
inundation variables.

\subsection{Data Processing and Transformation}
\label{sec:res-dat-cas-data-processing-transformation}

The integrated data workflow is:

\begin{enumerate}
  \item register source assets in the dataset manifest;
  \item choose final $P_{\mathrm{e}}$, Kamaishi target and Sendai
    target coordinates;
  \item recompute $\theta_{\mathrm{K}}$ and $\theta_{\mathrm{S}}$;
  \item generate corridor polygons and buffer/sponge polygons;
  \item clip and reconcile terrain, source and observation datasets;
  \item interpolate terrain and bed increments to regional mesh
    fields;
  \item extract regional interface histories for local 3D replay;
  \item write structured HDF5 or equivalent outputs with provenance.
\end{enumerate}

\subsection{Uncertainty, Provenance and Licensing}
\label{sec:res-dat-cas-uncertainty-provenance-licensing}

The integrated case shall not be accepted unless every source in
\cref{tab:res-dat-cas-data-source-decision-table} has a provenance
record, licence/usage note and downstream role. Missing-source records
are allowed only as explicit blockers or TODOs; they shall not be
converted into fabricated citation keys.

\subsection{Data Acceptance Criteria}
\label{sec:res-dat-cas-data-acceptance-criteria}

An impact case is accepted only after event reconstruction validation and
defence-model credibility are established. The case record shall identify the
event baseline, local exposure, defence geometry version, material and
foundation assumptions, numerical model fidelity, convergence state, physical
validity state and uncertainty state.

For the current hydrodynamic baseline, an integrated case is accepted
for software implementation planning when it contains the dataset
manifest, coordinate transform, corridor polygons, clipped terrain,
source candidate records, observation bundles and local replay-boundary
format. This acceptance does not mean validation has been executed.

\subsection{Candidate Approaches}
\label{sec:res-dat-cas-candidate-approaches}

% TODO: Define the credible candidate approaches considered.

\subsection{Critical Comparison}
\label{sec:res-dat-cas-critical-comparison}

% TODO: Compare candidates using physically and project-relevant criteria.
% Do not select an approach solely on popularity or implementation ease.

\subsection{Assumptions and Applicability}
\label{sec:res-dat-cas-assumptions}

% TODO: State assumptions and the conditions under which they are valid.

\subsection{Limitations and Failure Conditions}
\label{sec:res-dat-cas-limitations}

% TODO: State where the evidence, model or selected approach ceases to be
% reliable or applicable.

\subsection{Project-Specific Conclusions}
\label{sec:res-dat-cas-conclusions}

The initial case-study specification is event-conditioned and
defence-specific. Secondary events and cross-event holdouts are deferred until
the single-event workflow is complete and reviewed.

The current integrated case is not a generic literature review. Each
source is retained because it supports, limits or defers a project
decision: terrain construction, source forcing, observation validation,
corridor definition, optional sensitivity or missing defence geometry.

\subsection{Selected Approach and Rationale}
\label{sec:res-dat-cas-selected-approach}

Use a manifest-driven Tōhoku case-study specification:

\begin{quote}
  USGS event identity and source candidates
  $+$
  J-EGG500/GSI/GEBCO/ETOPO terrain hierarchy
  $+$
  NOAA/NCEI and Mori/JTS observations
  $+$
  Kamaishi and Sendai corridor polygons
  $+$
  structured local replay-boundary outputs
\end{quote}

ERA5 and JRA-55 are explicitly optional sensitivity sources rather
than baseline forcing. Kamaishi defence records remain a missing-source
blocker for real defence recreation.

\subsection{Unresolved Decisions}
\label{sec:res-dat-cas-unresolved-decisions}

Open decisions are final GIS target coordinates, final terrain product
licence acceptance, final finite-fault/source selection, final
calibration/validation split, tide-gauge station selection, Kamaishi
defence-source acquisition, final corridor dimensions and final replay
boundary schema.

\subsection{Requirements and Handoffs}
\label{sec:res-dat-cas-requirements}

Data handling now begins in parallel with remaining research. Software
shall start with a modular MVP containing:

\begin{itemize}
  \item finite-volume data model and mesh/field containers;
  \item manifest-driven raster/vector I/O;
  \item coordinate-transform and corridor-polygon generation;
  \item bathymetry/topography interpolation and clipping;
  \item regional 2D scaffold and boundary framework;
  \item HDF5 or equivalent structured output schema;
  \item replay-boundary format for local 3D inlet forcing;
  \item local 3D scaffold that consumes replay histories without
    implementing the solver in this Research deliverable.
\end{itemize}

`RES-VER-SPEC` shall consume the same manifest when deciding which
records are calibration, validation or diagnostic evidence.

\subsection{Deliverable Completion Record}
\label{sec:res-dat-cas-completion-record}

% TODO: Record whether the Deliverable is:
%
% - Not started
% - In progress
% - Draft complete
% - Reviewed
% - Accepted
%
% Acceptance requires:
%
% - the research questions to be answered;
% - claims to be supported by appropriate evidence;
% - assumptions and limitations to be explicit;
% - the project-specific conclusion to be stated;
% - downstream requirements to be traceable;
% - unresolved decisions to be closed or formally deferred.
